\documentclass[11pt,a4paper]{article}
\usepackage[utf8]{inputenc}
\usepackage{amsmath,amssymb,amsthm}
\usepackage{algorithm}
\usepackage{algpseudocode}
\usepackage{graphicx}
\usepackage{hyperref}
\usepackage{booktabs}
\usepackage{bm}
\usepackage{bbm}

% Custom commands
\newcommand{\R}{\mathbb{R}}
\newcommand{\E}{\mathbb{E}}
\newcommand{\prob}{\mathbb{P}}
\newcommand{\vx}{\mathbf{x}}
\newcommand{\vz}{\mathbf{z}}
\newcommand{\vtheta}{\boldsymbol{\theta}}
\newcommand{\sigmoid}{\sigma}
\DeclareMathOperator{\softmax}{softmax}

\title{Discrete Denoising Diffusion EDA with Gumbel-Softmax:\\
A Technical Documentation}
\author{Pateda Framework}
\date{\today}

\begin{document}

\maketitle

\begin{abstract}
This document provides a comprehensive mathematical description of the Discrete Denoising Diffusion Estimation of Distribution Algorithm (Dendiff-EDA) with Gumbel-Softmax. The algorithm combines principles from denoising diffusion probabilistic models with estimation of distribution algorithms for solving binary optimization problems. We describe the forward diffusion process, the neural network architecture for denoising, the learning procedure including various loss functions, the sampling mechanism, and the influence of key hyperparameters on algorithm behavior.
\end{abstract}

\tableofcontents
\newpage

%==============================================================================
\section{Introduction}
%==============================================================================

Estimation of Distribution Algorithms (EDAs) are population-based optimization methods that learn and sample from probabilistic models of promising solutions. The Discrete Denoising Diffusion EDA (Dendiff-EDA) introduces a novel approach by leveraging denoising diffusion probabilistic models (DDPMs) adapted for binary optimization.

The key innovations of Dendiff-Gumbel include:
\begin{enumerate}
    \item \textbf{Discrete diffusion process}: Modeling noise as bit-flip probability rather than Gaussian noise
    \item \textbf{Gumbel-Softmax reparameterization}: Enabling differentiable discrete sampling during training
    \item \textbf{Probabilistic denoising}: Learning to predict original binary values from corrupted inputs
    \item \textbf{Fitness-guided generation}: Optional conditioning on fitness values to bias sampling toward high-quality solutions
\end{enumerate}

%==============================================================================
\section{Problem Setting}
%==============================================================================

Consider a binary optimization problem where we seek to maximize (or minimize) an objective function $f: \{0,1\}^n \rightarrow \R$ over $n$ binary variables. Given a population of $N$ selected individuals $\mathcal{P} = \{\vx^{(1)}, \ldots, \vx^{(N)}\}$ where each $\vx^{(i)} \in \{0,1\}^n$, and their associated fitness values $\{f^{(1)}, \ldots, f^{(N)}\}$, the goal is to learn a generative model that can sample new candidate solutions likely to have high fitness.

%==============================================================================
\section{Forward Diffusion Process}
%==============================================================================

\subsection{Noise Schedule}

The forward diffusion process progressively corrupts binary data over $T$ timesteps. Unlike continuous diffusion models that add Gaussian noise, the discrete diffusion process corrupts data by randomly flipping bits.

Let $\beta_t$ denote the bit-flip probability at timestep $t$, where $t \in \{0, 1, \ldots, T-1\}$. The noise schedule defines how $\beta_t$ varies across timesteps:

\subsubsection{Linear Schedule}
\begin{equation}
    \beta_t = \beta_{\text{start}} + \frac{t}{T-1}(\beta_{\text{end}} - \beta_{\text{start}})
\end{equation}

\subsubsection{Cosine Schedule}
\begin{equation}
    \bar{\alpha}_t = \cos^2\left(\frac{t/T + 0.008}{1.008} \cdot \frac{\pi}{2}\right)
\end{equation}
\begin{equation}
    \beta_t = \text{clip}\left(1 - \frac{\bar{\alpha}_t}{\bar{\alpha}_{t-1}}, \beta_{\text{start}}, \beta_{\text{end}}\right)
\end{equation}

\subsection{Cumulative Noise Parameters}

Define the ``survival probability'' at each timestep:
\begin{equation}
    \alpha_t = 1 - \beta_t
\end{equation}

The cumulative product represents the probability that a bit remains unchanged after $t$ steps:
\begin{equation}
    \bar{\alpha}_t = \prod_{s=0}^{t} \alpha_s
\end{equation}

\subsection{Forward Corruption Process}

Given clean binary data $\vx_0 \in \{0,1\}^n$, the corrupted data at timestep $t$ is obtained via:
\begin{equation}
    q(\vx_t | \vx_0) = \text{Bernoulli-Flip}(\vx_0, 1 - \bar{\alpha}_t)
\end{equation}

Specifically, for each variable $i$:
\begin{equation}
    x_{t,i} = \begin{cases}
        1 - x_{0,i} & \text{with probability } 1 - \bar{\alpha}_t \\
        x_{0,i} & \text{with probability } \bar{\alpha}_t
    \end{cases}
\end{equation}

This can be implemented as:
\begin{equation}
    x_{t,i} = (x_{0,i} + \mathbbm{1}[u_i > \bar{\alpha}_t]) \mod 2
\end{equation}
where $u_i \sim \text{Uniform}(0,1)$.

%==============================================================================
\section{Neural Network Architecture}
%==============================================================================

\subsection{Denoising Network Structure}

The core component is a Multi-Layer Perceptron (MLP) that learns to predict the original binary values given corrupted input and the timestep. The network architecture consists of:

\begin{enumerate}
    \item \textbf{Input layer}: Concatenation of corrupted binary data $\vx_t \in \R^n$ with time embedding $\mathbf{e}_t \in \R^{d_t}$
    \item \textbf{Hidden layers}: Sequence of fully-connected layers with activation functions and dropout
    \item \textbf{Output layer}: Produces logits for binary classification $\mathbf{z} \in \R^{n \times 2}$
\end{enumerate}

\subsubsection{Time Embedding}

The timestep $t$ is encoded using sinusoidal positional embeddings:
\begin{equation}
    \mathbf{e}_t[2k] = \sin\left(\frac{t}{10000^{2k/d_t}}\right)
\end{equation}
\begin{equation}
    \mathbf{e}_t[2k+1] = \cos\left(\frac{t}{10000^{2k/d_t}}\right)
\end{equation}
where $d_t$ is the time embedding dimension and $k \in \{0, 1, \ldots, d_t/2 - 1\}$.

\subsubsection{Network Forward Pass}

Let $\vtheta$ denote the network parameters. The denoising network computes:
\begin{equation}
    \mathbf{h}_0 = [\vx_t; \mathbf{e}_t] \in \R^{n + d_t}
\end{equation}
\begin{equation}
    \mathbf{h}_\ell = \text{Dropout}\left(\phi_\ell\left(W_\ell \mathbf{h}_{\ell-1} + \mathbf{b}_\ell\right)\right), \quad \ell = 1, \ldots, L
\end{equation}
\begin{equation}
    \mathbf{z} = W_{\text{out}} \mathbf{h}_L + \mathbf{b}_{\text{out}} \in \R^{n \times 2}
\end{equation}

where $\phi_\ell$ is the activation function (e.g., ReLU, ELU, GELU) for layer $\ell$.

\subsection{Network Inputs and Outputs}

\begin{table}[h]
\centering
\begin{tabular}{@{}lll@{}}
\toprule
\textbf{Component} & \textbf{Shape} & \textbf{Description} \\
\midrule
\textbf{Inputs:} & & \\
$\vx_t$ & $(B, n)$ & Corrupted binary data \\
$t$ & $(B,)$ & Timestep indices \\
$\mathbf{e}_t$ & $(B, d_t)$ & Time embeddings \\
$f$ (optional) & $(B, 1)$ & Fitness values (for guided version) \\
\midrule
\textbf{Outputs:} & & \\
$\mathbf{z}$ & $(B, n, 2)$ & Logits for binary classification \\
$p(\vx_0 | \vx_t, t)$ & $(B, n, 2)$ & Predicted probabilities after softmax \\
\bottomrule
\end{tabular}
\caption{Network inputs and outputs, where $B$ is batch size, $n$ is number of variables, and $d_t$ is time embedding dimension.}
\label{tab:io}
\end{table}

\subsection{Probability Interpretation}

The output logits are interpreted as:
\begin{equation}
    z_{i,0} = \log p(x_{0,i} = 0 | \vx_t, t), \quad z_{i,1} = \log p(x_{0,i} = 1 | \vx_t, t)
\end{equation}

Applying softmax yields the predicted probability distribution:
\begin{equation}
    p(x_{0,i} = c | \vx_t, t) = \frac{\exp(z_{i,c})}{\exp(z_{i,0}) + \exp(z_{i,1})}, \quad c \in \{0, 1\}
\end{equation}

%==============================================================================
\section{Gumbel-Softmax Reparameterization}
%==============================================================================

\subsection{The Discrete Sampling Challenge}

Training requires backpropagation through discrete samples, which is non-differentiable. The Gumbel-Softmax trick provides a continuous relaxation.

\subsection{Gumbel-Softmax Distribution}

Given logits $\mathbf{z} \in \R^{n \times 2}$ and temperature $\tau > 0$:

\begin{enumerate}
    \item Sample Gumbel noise:
    \begin{equation}
        g_{i,c} = -\log(-\log(u_{i,c})), \quad u_{i,c} \sim \text{Uniform}(0,1)
    \end{equation}

    \item Compute soft samples:
    \begin{equation}
        y_{i,c} = \frac{\exp((z_{i,c} + g_{i,c})/\tau)}{\sum_{c'=0}^{1} \exp((z_{i,c'} + g_{i,c'})/\tau)}
    \end{equation}
\end{enumerate}

\subsection{Straight-Through Estimator}

For hard samples with gradient flow, use the straight-through estimator:
\begin{equation}
    \hat{y}_{i,c} = \begin{cases}
        1 & \text{if } c = \arg\max_{c'} y_{i,c'} \\
        0 & \text{otherwise}
    \end{cases}
\end{equation}
\begin{equation}
    \tilde{y} = \hat{y} - \text{sg}(y) + y
\end{equation}
where $\text{sg}(\cdot)$ is the stop-gradient operator, allowing gradients to flow through $y$ while using $\hat{y}$ in the forward pass.

\subsection{Temperature Annealing}

The temperature controls the ``hardness'' of the samples:
\begin{itemize}
    \item $\tau \rightarrow 0$: Samples approach one-hot (discrete)
    \item $\tau \rightarrow \infty$: Samples approach uniform (soft)
\end{itemize}

Temperature is typically annealed during training:
\begin{equation}
    \tau_{e+1} = \max(\tau_{\min}, \tau_e \cdot \gamma)
\end{equation}
where $\gamma$ is the decay rate (e.g., 0.99) and $\tau_{\min}$ is the minimum temperature (e.g., 0.5).

%==============================================================================
\section{Learning Algorithm}
%==============================================================================

\subsection{Training Objective}

The model is trained to predict the original binary values from corrupted inputs. The primary objective is cross-entropy loss:
\begin{equation}
    \mathcal{L}_{\text{CE}}(\vtheta) = -\E_{t, \vx_0, \vx_t}\left[\sum_{i=1}^{n} \log p_\vtheta(x_{0,i} | \vx_t, t)\right]
\end{equation}

where the expectation is over:
\begin{itemize}
    \item Timesteps $t \sim \text{Uniform}(\{0, \ldots, T-1\})$
    \item Clean data $\vx_0 \sim p_{\text{data}}$ (the selected population)
    \item Corrupted data $\vx_t \sim q(\vx_t | \vx_0)$
\end{itemize}

\subsection{Learning Procedure}

The \texttt{learn\_discrete\_dendiff\_gumbel\_enhanced()} function implements the training procedure:

\begin{algorithm}[H]
\caption{Learning Discrete Dendiff-Gumbel Model}
\label{alg:learning}
\begin{algorithmic}[1]
\Require Population $\mathcal{P} = \{\vx^{(1)}, \ldots, \vx^{(N)}\}$, fitness values $\{f^{(1)}, \ldots, f^{(N)}\}$
\Require Parameters: $T$ (timesteps), $\beta_{\text{start}}$, $\beta_{\text{end}}$, epochs $E$, batch size $B$
\State Initialize denoising network $D_\vtheta$ with parameters $\vtheta$
\State Compute beta schedule: $\beta_0, \ldots, \beta_{T-1}$
\State Compute cumulative alphas: $\bar{\alpha}_0, \ldots, \bar{\alpha}_{T-1}$
\State Set initial temperature $\tau \gets \tau_{\text{init}}$
\For{epoch $= 1$ to $E$}
    \State Shuffle dataset indices
    \For{each mini-batch $\mathcal{B} \subset \mathcal{P}$}
        \State Sample timesteps: $t \sim \text{Uniform}(\{0, \ldots, T-1\})$ for each sample
        \State Corrupt data: $\vx_t \gets q(\vx_t | \vx_0, t)$ \Comment{Forward diffusion}
        \State Forward pass: $\mathbf{z} \gets D_\vtheta(\vx_t, t)$ \Comment{Get logits}
        \State Compute loss: $\mathcal{L} \gets \text{Loss}(\mathbf{z}, \vx_0, \mathbf{f})$
        \State Backward pass: $\nabla_\vtheta \mathcal{L}$
        \State Clip gradients: $\|\nabla_\vtheta\| \gets \min(\|\nabla_\vtheta\|, 1.0)$
        \State Update parameters: $\vtheta \gets \vtheta - \eta \nabla_\vtheta \mathcal{L}$
    \EndFor
    \State Anneal temperature: $\tau \gets \max(\tau_{\min}, \tau \cdot \gamma)$
\EndFor
\State \Return trained model $D_\vtheta$
\end{algorithmic}
\end{algorithm}

%==============================================================================
\section{Loss Functions}
%==============================================================================

The implementation supports multiple loss functions, each with different properties for guiding the learning process.

\subsection{Standard Cross-Entropy Loss (MSE)}

The default loss function treats the problem as classification:
\begin{equation}
    \mathcal{L}_{\text{CE}} = -\frac{1}{NB}\sum_{j=1}^{B}\sum_{i=1}^{n} \left[x_{0,i}^{(j)} \log \hat{p}_{i,1}^{(j)} + (1-x_{0,i}^{(j)}) \log \hat{p}_{i,0}^{(j)}\right]
\end{equation}
where $\hat{p}_{i,c}^{(j)} = \softmax(\mathbf{z}^{(j)})_{i,c}$ is the predicted probability.

\subsection{Weighted Cross-Entropy Loss (Weighted MSE)}

Higher fitness samples receive higher weight in the loss:
\begin{equation}
    w^{(j)} = \frac{f^{(j)} - f_{\min}}{f_{\max} - f_{\min} + \epsilon}
\end{equation}
\begin{equation}
    \mathcal{L}_{\text{weighted}} = \frac{1}{\sum_j w^{(j)}}\sum_{j=1}^{B} w^{(j)} \cdot \mathcal{L}_{\text{CE}}^{(j)}
\end{equation}

This encourages the model to more accurately reconstruct high-fitness solutions.

\subsection{Ranking Loss}

The ranking loss prioritizes learning the relative ordering of solutions:
\begin{equation}
    \mathcal{L}_{\text{rank}} = (1-\lambda) \cdot \mathcal{L}_{\text{CE}} + \lambda \cdot \mathcal{L}_{\text{pairwise}}
\end{equation}

where the pairwise ranking component is:
\begin{equation}
    \mathcal{L}_{\text{pairwise}} = \frac{1}{|\mathcal{S}|} \sum_{(j,k) \in \mathcal{S}} \max\left(0, \ell^{(j)} - \ell^{(k)} + m\right)
\end{equation}

where $\mathcal{S} = \{(j,k) : f^{(j)} > f^{(k)}\}$ is the set of valid pairs, $\ell^{(j)}$ is the per-sample loss for sample $j$, and $m$ is the margin parameter.

This loss encourages the model to produce lower reconstruction error for higher-fitness samples.

\subsection{Huber Loss}

The Huber loss provides robustness to outliers:
\begin{equation}
    L_\delta(a) = \begin{cases}
        \frac{1}{2}a^2 & \text{if } |a| \leq \delta \\
        \delta(|a| - \frac{1}{2}\delta) & \text{otherwise}
    \end{cases}
\end{equation}

Applied to negative log-likelihood:
\begin{equation}
    \mathcal{L}_{\text{Huber}} = \frac{1}{Bn}\sum_{j=1}^{B}\sum_{i=1}^{n} L_\delta\left(-\log \hat{p}_{i,x_{0,i}}^{(j)}\right)
\end{equation}

%==============================================================================
\section{Fitness-Guided Generation}
%==============================================================================

\subsection{Conditional Denoising Network}

The fitness-guided variant conditions the denoising on fitness information, inspired by conditional VAEs (C-VAE) and fitness-guided approaches.

The network architecture is modified to include fitness embedding:
\begin{equation}
    \mathbf{e}_f = W_f \cdot \tilde{f} + \mathbf{b}_f \in \R^{d_f}
\end{equation}
where $\tilde{f} = (f - f_{\min})/(f_{\max} - f_{\min})$ is the normalized fitness.

The input to the hidden layers becomes:
\begin{equation}
    \mathbf{h}_0 = [\vx_t; \mathbf{e}_t; \mathbf{e}_f] \in \R^{n + d_t + d_f}
\end{equation}

\subsection{Biasing Mechanism}

During training, the network learns the conditional distribution:
\begin{equation}
    p_\vtheta(\vx_0 | \vx_t, t, f)
\end{equation}

This enables the network to:
\begin{enumerate}
    \item Learn that high-fitness solutions have certain structural patterns
    \item Condition the denoising process on desired fitness levels
    \item Generate solutions biased toward high-fitness regions of the search space
\end{enumerate}

\subsection{Guided Sampling}

During sampling, we set a target fitness value (typically the maximum normalized fitness $\tilde{f} = 1.0$):
\begin{equation}
    \vx_0 \sim p_\vtheta(\cdot | \vx_T, T \rightarrow 0, \tilde{f}_{\text{target}})
\end{equation}

This biases the generation toward solutions that the model associates with high fitness.

%==============================================================================
\section{Sampling Algorithm}
%==============================================================================

\subsection{Reverse Diffusion Process}

The sampling procedure iteratively denoises random binary data:

\begin{algorithm}[H]
\caption{Sampling from Discrete Dendiff-Gumbel}
\label{alg:sampling}
\begin{algorithmic}[1]
\Require Trained model $D_\vtheta$, number of samples $M$, sampling steps $S$, temperature $\tau$
\State Initialize: $\vx_T \sim \text{Bernoulli}(0.5)^{M \times n}$ \Comment{Random binary}
\If{fitness-guided}
    \State Set target fitness: $\tilde{f}_{\text{target}} \gets 1.0$ \Comment{Aim for high fitness}
\EndIf
\State Create timestep schedule: $\{t_1, t_2, \ldots, t_S\}$ from $T-1$ to $0$
\For{$s = 1$ to $S$}
    \State $t \gets t_s$
    \If{fitness-guided}
        \State $\mathbf{z} \gets D_\vtheta(\vx_t, t, \tilde{f}_{\text{target}})$
    \Else
        \State $\mathbf{z} \gets D_\vtheta(\vx_t, t)$
    \EndIf
    \If{deterministic}
        \State $\hat{\vx}_0 \gets \arg\max_c \mathbf{z}_{:,:,c}$
    \Else
        \State $\mathbf{y} \gets \text{Gumbel-Softmax}(\mathbf{z}, \tau, \text{hard}=\texttt{True})$
        \State $\hat{\vx}_0 \gets \mathbf{y}_{:,:,1}$ \Comment{Probability of bit being 1}
    \EndIf
    \If{$t > 0$}
        \State $\rho \gets \text{clip}(\bar{\alpha}_{t-1} / \bar{\alpha}_t, 0, 1)$ \Comment{Mixing probability}
        \State $\mathbf{m} \gets (\mathbf{u} < \rho)$ where $\mathbf{u} \sim \text{Uniform}(0,1)^{M \times n}$
        \State $\vx_{t-1} \gets \mathbf{m} \odot \hat{\vx}_0 + (1 - \mathbf{m}) \odot \vx_t$
    \Else
        \State $\vx_0 \gets \hat{\vx}_0$
    \EndIf
\EndFor
\State \Return $(\vx_0 > 0.5)$ \Comment{Binarize final output}
\end{algorithmic}
\end{algorithm}

\subsection{Progressive Denoising}

The key insight is the progressive mixing between predicted clean data and current noisy data:
\begin{equation}
    \vx_{t-1} = \rho_t \cdot \hat{\vx}_0 + (1 - \rho_t) \cdot \vx_t
\end{equation}
where $\rho_t = \bar{\alpha}_{t-1}/\bar{\alpha}_t$ represents the trust in the prediction.

As $t$ decreases:
\begin{itemize}
    \item $\rho_t$ increases (more trust in prediction)
    \item The corruption level decreases
    \item The samples converge to the learned distribution
\end{itemize}

\subsection{Strided Sampling}

For faster sampling with $S < T$ steps, a strided timestep schedule is used:
\begin{equation}
    t_s = \left\lfloor (T-1) \cdot \frac{S - s}{S - 1} \right\rfloor, \quad s = 1, \ldots, S
\end{equation}

%==============================================================================
\section{Complete EDA Algorithm}
%==============================================================================

\begin{algorithm}[H]
\caption{Discrete Dendiff-Gumbel EDA}
\label{alg:eda}
\begin{algorithmic}[1]
\Require Fitness function $f$, population size $N$, generations $G$, selection ratio $\rho$
\Require Dendiff parameters: $T$, $\beta_{\text{start}}$, $\beta_{\text{end}}$, epochs, $\tau$
\State Initialize population: $\mathcal{P}_0 \sim \text{Bernoulli}(0.5)^{N \times n}$
\State Evaluate: $f^{(i)} \gets f(\vx^{(i)})$ for all $i$
\State Track best: $\vx^* \gets \arg\max_{\vx \in \mathcal{P}_0} f(\vx)$
\For{generation $g = 1$ to $G$}
    \State \textbf{Selection:}
    \State $k \gets \lfloor \rho \cdot N \rfloor$ \Comment{Number to select}
    \State $\mathcal{S}_g \gets $ top-$k$ individuals from $\mathcal{P}_{g-1}$ by fitness
    \State
    \State \textbf{Learning:}
    \State $D_\vtheta \gets \texttt{learn\_discrete\_dendiff\_gumbel\_enhanced}(\mathcal{S}_g, \mathbf{f}_{\mathcal{S}_g}, \text{params})$
    \State
    \State \textbf{Sampling:}
    \State $\mathcal{P}_g \gets \texttt{sample\_discrete\_dendiff\_gumbel}(D_\vtheta, N, \text{sampling\_params})$
    \State
    \State \textbf{Evaluation:}
    \State $f^{(i)} \gets f(\vx^{(i)})$ for all $\vx^{(i)} \in \mathcal{P}_g$
    \State
    \State \textbf{Update best:}
    \If{$\max_{\vx \in \mathcal{P}_g} f(\vx) > f(\vx^*)$}
        \State $\vx^* \gets \arg\max_{\vx \in \mathcal{P}_g} f(\vx)$
    \EndIf
    \State
    \State \textbf{Early termination:}
    \If{$|f(\vx^*) - f_{\text{optimal}}| < \epsilon$}
        \State \textbf{break}
    \EndIf
\EndFor
\State \Return $\vx^*$, $f(\vx^*)$
\end{algorithmic}
\end{algorithm}

%==============================================================================
\section{Hyperparameter Analysis}
%==============================================================================

\subsection{Temperature ($\tau$)}

The temperature parameter controls the sharpness of the Gumbel-Softmax distribution:

\begin{equation}
    y_{i,c} = \frac{\exp((z_{i,c} + g_{i,c})/\tau)}{\sum_{c'} \exp((z_{i,c'} + g_{i,c'})/\tau)}
\end{equation}

\textbf{Effects:}
\begin{itemize}
    \item $\tau \rightarrow 0$: Samples become one-hot (discrete), but gradients become unstable
    \item $\tau \rightarrow \infty$: Samples become uniform, losing discrimination
    \item $\tau \approx 0.5$--$1.0$: Good balance for training
\end{itemize}

\textbf{During training:} Temperature annealing (starting high, decreasing) helps:
\begin{enumerate}
    \item Early training: Softer samples enable smoother gradients
    \item Late training: Harder samples approximate true discrete behavior
\end{enumerate}

\textbf{During sampling:} Lower temperatures produce more deterministic, peaked distributions.

\subsection{Beta Start ($\beta_{\text{start}}$)}

The initial noise level in the diffusion schedule.

\textbf{Effects:}
\begin{itemize}
    \item \textbf{Low $\beta_{\text{start}}$ (e.g., 0.0001):}
        \begin{itemize}
            \item Early timesteps have minimal corruption
            \item Fine-grained denoising at early steps
            \item Better preservation of structure in early diffusion
        \end{itemize}
    \item \textbf{High $\beta_{\text{start}}$ (e.g., 0.01):}
        \begin{itemize}
            \item Faster initial corruption
            \item Coarser diffusion process
            \item May miss subtle patterns
        \end{itemize}
\end{itemize}

\textbf{Typical values:} $10^{-4}$ to $10^{-2}$

\subsection{Beta End ($\beta_{\text{end}}$)}

The final noise level in the diffusion schedule.

\textbf{Effects:}
\begin{itemize}
    \item \textbf{Low $\beta_{\text{end}}$ (e.g., 0.1):}
        \begin{itemize}
            \item Even at $t=T-1$, significant structure remains
            \item Easier denoising task
            \item May not fully explore the space
        \end{itemize}
    \item \textbf{High $\beta_{\text{end}}$ (e.g., 0.5):}
        \begin{itemize}
            \item Maximum corruption approaches random (50\% flip probability)
            \item Harder denoising task
            \item Better exploration but harder optimization
        \end{itemize}
\end{itemize}

\textbf{Important constraint:} $\beta_{\text{end}} < 0.5$ to maintain meaningful signal.

\textbf{Typical values:} $0.1$ to $0.4$

\subsection{Cumulative Effect}

The cumulative survival probability $\bar{\alpha}_T$ determines how much of the original signal remains:
\begin{equation}
    \bar{\alpha}_T = \prod_{t=0}^{T-1}(1 - \beta_t)
\end{equation}

For the linear schedule:
\begin{equation}
    \bar{\alpha}_T \approx \exp\left(-\frac{T}{2}(\beta_{\text{start}} + \beta_{\text{end}})\right)
\end{equation}

\textbf{Design principle:} Choose parameters such that $\bar{\alpha}_T \approx 0.01$--$0.1$, meaning 90-99\% of information is corrupted at the final timestep.

\subsection{Interaction with Timesteps ($T$)}

\begin{table}[h]
\centering
\begin{tabular}{@{}cccc@{}}
\toprule
$T$ & $\beta_{\text{start}}$ & $\beta_{\text{end}}$ & Behavior \\
\midrule
50 & 0.0001 & 0.3 & Moderate corruption, fast training \\
100 & 0.0001 & 0.3 & Finer diffusion, better quality \\
100 & 0.01 & 0.5 & Aggressive corruption, exploration \\
200 & 0.0001 & 0.2 & Very fine diffusion, slow training \\
\bottomrule
\end{tabular}
\caption{Example parameter configurations and their effects.}
\label{tab:params}
\end{table}

%==============================================================================
\section{Summary of Mathematical Formulations}
%==============================================================================

\subsection{Forward Process}
\begin{equation}
    q(\vx_t | \vx_0) = \prod_{i=1}^{n} \left[\bar{\alpha}_t \cdot \mathbbm{1}[x_{t,i} = x_{0,i}] + (1-\bar{\alpha}_t) \cdot \mathbbm{1}[x_{t,i} \neq x_{0,i}]\right]
\end{equation}

\subsection{Denoising Network}
\begin{equation}
    p_\vtheta(\vx_0 | \vx_t, t) = \prod_{i=1}^{n} \softmax(D_\vtheta(\vx_t, t))_i
\end{equation}

\subsection{Gumbel-Softmax}
\begin{equation}
    y = \softmax\left(\frac{\mathbf{z} + \mathbf{g}}{\tau}\right), \quad \mathbf{g} \sim \text{Gumbel}(0,1)
\end{equation}

\subsection{Loss Functions}
\begin{align}
    \mathcal{L}_{\text{CE}} &= -\E\left[\sum_i x_{0,i}\log\hat{p}_{i,1} + (1-x_{0,i})\log\hat{p}_{i,0}\right] \\
    \mathcal{L}_{\text{weighted}} &= \E\left[w(f) \cdot \mathcal{L}_{\text{CE}}\right] \\
    \mathcal{L}_{\text{ranking}} &= (1-\lambda)\mathcal{L}_{\text{CE}} + \lambda \cdot \E_{f_i > f_j}\left[\max(0, \ell_i - \ell_j + m)\right] \\
    \mathcal{L}_{\text{Huber}} &= \E\left[L_\delta(-\log \hat{p}_{x_0})\right]
\end{align}

\subsection{Reverse Process (Sampling)}
\begin{equation}
    \vx_{t-1} = \rho_t \cdot D_\vtheta(\vx_t, t) + (1-\rho_t) \cdot \vx_t, \quad \rho_t = \frac{\bar{\alpha}_{t-1}}{\bar{\alpha}_t}
\end{equation}

%==============================================================================
\section{Conclusion}
%==============================================================================

The Discrete Denoising Diffusion EDA with Gumbel-Softmax provides a powerful framework for binary optimization by:

\begin{enumerate}
    \item Leveraging diffusion models' ability to capture complex data distributions
    \item Using Gumbel-Softmax for differentiable discrete sampling
    \item Supporting multiple loss functions for different optimization scenarios
    \item Enabling fitness-guided generation to bias toward high-quality solutions
    \item Providing flexible hyperparameters ($\tau$, $\beta_{\text{start}}$, $\beta_{\text{end}}$) to control the diffusion-denoising process
\end{enumerate}

The algorithm is particularly suited for problems where:
\begin{itemize}
    \item Variable interactions are complex and difficult to model explicitly
    \item The fitness landscape has multiple local optima
    \item Exploration-exploitation balance is critical
\end{itemize}

\end{document}
